\section{Introduction}
\label{sec:intro}

AI compute is valuable but opaque. A customer pays a provider to run a specific
model on a specific input and ordinarily receives only an output and invoice.
The provider cannot turn the same execution into portable public evidence, and a
third party cannot independently distinguish the claimed workload from a copied
answer or fabricated meter.

Synthetic ``AI-like'' proof of work does not solve this problem. Replacing hash
grinding with a chain-derived matrix race demonstrates accelerator expenditure,
but the product is not necessarily requested or consumed by anyone. Requiring
every full node to replay that matrix work on a GPU also moves consensus toward a
specialized hardware requirement; delegating replay to signed GPU attestors
introduces a new trust authority.

Hanzo Proof of AI takes the opposite approach. The work is useful first: a real
customer or governed network service admits a job and fixes its consumer,
deliverable, model, input commitment, deterministic execution policy, proof
policy, reward, and payout destination before a miner returns a result. The
provider then commits evidence produced alongside that same execution. A-Chain
PoAI consensus determines that the bound work was actually done, consumes a
global nullifier, and finalizes one mining receipt.

\subsection{The cloud-provider proposition}

An existing cloud does not need to abandon its scheduler or customer contract.
It can attach the PoAI adapter to inference, embedding, evaluation, or supported
training work it already runs. The cloud retains normal service revenue and may
earn declining AI mining rewards for making that execution independently
verifiable. A decentralized provider or idle workstation accepts open jobs
through the same protocol.

This is a ``compute once, earn twice'' market, not free issuance. Proof capture
must be materially cheaper than the original workload; the job must be admitted
independently of the miner; metering must be verifiable; and the same work
authorization must never be rewarded twice.

\subsection{CPU-verifiable work}

The canonical public path uses deterministic integer execution so honest CPU and
GPU backends commit the same bytes. Small jobs may be replayed exactly. For a
large committed matrix product, an ordinary CPU can verify a post-commitment
Freivalds opening in quadratic work rather than repeating the cubic product.
Merkle commitments bind that opening into the model graph.

This arithmetic asymmetry is the core, not the complete protocol. Whole-job
accountability additionally requires graph wiring, unpredictable challenges,
coverage, transcript data availability, proof finality, and verifiable metering.
TEE evidence may be added for confidential jobs. An activated P3Q proof may
eventually replace optimistic challenges, but generic proof-system code is not
production PoAI authority.

\subsection{PoAI consensus and P3Q settlement on the Lux PQ base layer}

The three layers answer different questions:

\begin{itemize}[leftmargin=1.2em]
  \item \textbf{PoAI consensus on A-Chain} decides whether useful work was
        actually performed, whether its evidence is final, and whether it may
        authorize an AI reward.
  \item \textbf{P3Q settlement on Z-Chain} batches finalized A-Chain receipt
        transitions into a Plonky3-derived, pairing-free STARK/FRI proof. It
        compresses and settles A-Chain decisions; it cannot admit raw work or
        authorize mining.
  \item \textbf{Lux Quasar} orders and finalizes the A-Chain and Z-Chain
        transitions and supplies the activated post-quantum validator context
        for exporting their settlement root.
\end{itemize}

There is one mining authority. ``Mine into another chain'' means the A-Chain job
selects that chain as the payout destination. The destination verifies the
Z-settled, PQ-authenticated A-Chain receipt and consumes it exactly once; it
cannot accept raw work or run an independent emission schedule.

\subsection{The AI economy}

AI is the native coin of Hanzo Network. Its maximum nominal supply is two
trillion: one trillion permanently allocated to the Hanzo DAO and one trillion
reserved for proof-earned miner and validator rewards. The mining pool follows a
Bitcoin-shaped geometric schedule: 500 billion in the first era, 250 billion in
the second, 125 billion in the third, and so on. Fee burns may make net supply
deflationary as issuance declines.

\subsection{Contributions and outline}

This paper specifies:

\begin{enumerate}[leftmargin=1.4em]
  \item the A-Chain PoAI useful-work state machine, Z-Chain P3Q settlement, and
        their relation to Lux post-quantum finality (\S\ref{sec:arch});
  \item the two-trillion AI hard cap, fixed DAO allocation, one-trillion
        proof-earned pool, halving schedule, fees, staking, and burns
        (\S\ref{sec:aicoin});
  \item the HMM model registry and agent identity surfaces
        (\S\ref{sec:hmm}, \S\ref{sec:agents});
  \item proof-bearing inference, training, evaluation, and capacity markets,
        including the existing-cloud adapter (\S\ref{sec:markets});
  \item receipt-only omnichain settlement rooted in A-Chain, Z-Chain, and Lux
        (\S\ref{sec:xchain}); and
  \item confidentiality and security boundaries, with explicit launch gates
        rather than assumed deployment (\S\ref{sec:privacy},
        \S\ref{sec:security}, \S\ref{sec:deploy}).
\end{enumerate}
